% 1
\begin{frame}[plain]
  \titlepage
  \vspace{-1.2em}
  \begin{center}
  \small Dla inżynierów: jasno, ale bez ukrywania modelu, ograniczeń i ryzyk.
  \end{center}
  \note{Slajd ustawia cel prezentacji: nie jest to gotowa obietnica produktu, tylko początek projektu badawczego nad sterowaniem predykcyjnym i bezpieczeństwem lokomocji.}
\end{frame}

% 2
\begin{frame}{Cel projektu w jednym obrazie}
\begin{columns}[T,onlytextwidth]
\column{0.56\textwidth}
\begin{itemize}
  \item Mamy czworonożnego robota z nietypową nogą opartą o mechanizm czworoboku.
  \item Obecne sterowanie oparte na \key{Reinforcement Learning}, czyli uczeniu przez wzmacnianie, działa w wielu próbach.
  \item Problem: polityka może wprowadzić robota w konfiguracje trudne do odzyskania na realnym podłożu.
  \item Hipoteza projektu: \key{model + predykcja + bariery bezpieczeństwa} mogą zmniejszyć liczbę takich awarii.
\end{itemize}
\column{0.40\textwidth}
\begin{center}
\begin{tikzpicture}[scale=0.88,>=Latex]
  \draw[thick,rounded corners=3pt,fill=blue!8] (-1.7,0.25) rectangle (1.7,1.05);
  \node[font=\small] at (0,0.65) {korpus robota};
  \foreach \x/\d in {-1.25/-1,-0.45/1,0.45/-1,1.25/1} {
    \draw[very thick] (\x,0.25)--(\x+0.25*\d,-0.65)--(\x+0.55*\d,-1.35);
    \fill[safe] (\x+0.55*\d,-1.40) circle (0.07);
  }
  \draw[->,thick,pwrblue] (-1.9,1.4)--(1.9,1.4);
  \node[font=\scriptsize] at (0,1.65) {ruch korpusu i harmonogram kontaktów};
  \node[font=\scriptsize,align=center] at (0,-1.95) {RL proponuje ruch\\MPC i CBF kontrolują ryzyko};
\end{tikzpicture}
\end{center}

Czym jest MPC: przewidywanie, optymalizacja, ograniczenia, sterowanie w pętli.
\end{columns}
\source{Opis platformy z KKR: cztery kończyny, 3 stopnie swobody na kończynę, ROS 2, CAN, tryby sterowania napędów.}
\note{Należy wyjaśnić, że projekt nie neguje uczenia przez wzmacnianie. Chodzi o dodanie warstwy modelowej i bezpieczeństwa tam, gdzie czysta polityka uczona nie ma jawnych gwarancji.}
\end{frame}

% 3
\begin{frame}{Platforma: co wiemy z dokumentacji}
\begin{itemize}
  \item Korpus aluminiowy i cztery kończyny.
  \item Każda kończyna: \key{trzy stopnie swobody}.
  \item Główna część nogi: mechanizm czworoboku przegubowego / równoległoboku.
  \item Napędy: silniki bezszczotkowe prądu stałego z przekładniami planetarnymi.
  \item Sterowanie z komputera przez \key{Robot Operating System 2 (ROS 2)} i \key{Controller Area Network (CAN)}.
\end{itemize}
{\vspace{0.15em}\tiny\color{black!55}Źródło: załączony PDF, s. 1-2. Pełne dane CAD, bezwładności i ograniczenia momentu wymagają uzupełnienia.}
\note{Tu trzeba wyraźnie odróżnić fakty z PDF-a od przyszłych założeń modelowych. Dokładny model wymaga CAD-u, danych silników, przekładni i identyfikacji.}
\end{frame}

% 4
\begin{frame}{Dlaczego to nie jest „zwykły manipulator na nogach”?}
\begin{columns}[T,onlytextwidth]
\column{0.52\textwidth}
\begin{itemize}
  \item Robot kroczący zmienia punkty kontaktu z podłożem.
  \item Kontakt może się pojawić, zniknąć, poślizgnąć albo uderzyć w podłoże.
  \item Tarcie i nierówności podłoża są częścią sterowania, nie tylko zakłóceniem.
  \item Błąd jednej nogi może zmienić równowagę całego korpusu.
  \item Kontakt i podłoże są aktywnymi elementami układu.
\end{itemize}
\column{0.44\textwidth}
\begin{block}{Konsekwencja dla sterowania}
Sterownik musi planować ruch, siły i bezpieczeństwo konfiguracji, a nie tylko śledzić zadaną trajektorię przegubów.
\end{block}
\vspace{0.5em}
\begin{center}
\begin{tikzpicture}[scale=0.9]
  \draw[thick,rounded corners=2pt,fill=blue!6] (-1.5,0.2) rectangle (1.5,0.8);
  \foreach \x/\s in {-1.2/1, -0.45/0, 0.45/1, 1.2/0} {
    \draw[thick] (\x,0.2) -- (\x-0.25,-0.8);
    \draw[thick] (\x,0.2) -- (\x+0.25,-0.8);
    \ifnum\s=1 \fill[safe] (\x,-0.85) circle (0.08); \else \draw[warn,thick] (\x,-0.85) circle (0.08); \fi
  }
  \draw[black!40,thick] (-2,-0.95) -- (2,-0.95);
  \node[align=center,font=\scriptsize] at (0,-1.35) {kontakt aktywny / noga w przenoszeniu};
\end{tikzpicture}
\end{center}
\end{columns}
\note{To dobry slajd dla laików: kontakt i podłoże są aktywnymi elementami układu. W sterowaniu manipulatorów zwykle mamy stałą podstawę, a tutaj podstawa musi być utrzymywana przez same nogi.}
\end{frame}

% 5
\begin{frame}{Obecny punkt startowy: uczenie przez wzmacnianie}
\small
\begin{columns}[T,onlytextwidth]
\column{0.49\textwidth}
\begin{block}{Reinforcement Learning (RL)}
Polityka sterowania uczona przez próby i nagrody w symulacji lub na robocie.
\end{block}
\begin{itemize}
  \item Znajduje złożone zachowania bez ręcznego programowania każdego kroku.
  \item Po treningu działa szybko.
  \item Bezpieczeństwo bywa zapisane tylko pośrednio w funkcji nagrody.
\end{itemize}
\column{0.47\textwidth}
\begin{center}
\begin{tikzpicture}[node distance=0.65cm,>=Latex]
  \tikzstyle{box}=[draw,rounded corners=2pt,fill=blue!6,minimum width=3.0cm,minimum height=0.72cm,align=center,font=\small]
  \node[box] (obs) {obserwacje\\czujników};
  \node[box,below=of obs] (policy) {polityka RL};
  \node[box,below=of policy] (cmd) {zadane pozycje\\lub momenty};
  \node[box,below=of cmd] (robot) {robot + podłoże};
  \draw[->,thick] (obs) -- (policy);
  \draw[->,thick] (policy) -- (cmd);
  \draw[->,thick] (cmd) -- (robot);
  \draw[->,thick] (robot.west) .. controls +(-1.0,0) and +(-1.0,0) .. (obs.west);
\end{tikzpicture}
\end{center}
\end{columns}
\note{Należy podkreślić, że RL jest mocnym punktem startowym. Problemem nie jest brak ruchu, tylko brak jawnego mechanizmu wykluczania stanów niebezpiecznych lub nieodzyskiwalnych.}
\end{frame}

% 6
\begin{frame}{Awaria, która motywuje projekt}
\begin{columns}[T,onlytextwidth]
\column{0.52\textwidth}
\begin{itemize}
  \item Przykład: nogi ustawiają się zbyt szeroko na boki.
  \item Na szorstkim podłożu stopa nie przesuwa się łatwo do środka.
  \item Napęd może mieć za mały zapas momentu, aby odzyskać postawę.
  \item Polityka działała lokalnie sensownie, ale weszła w obszar o złej geometrii i złym kontakcie.
\end{itemize}
\column{0.44\textwidth}
\begin{block}{Nie chodzi tylko o „lepszy trening”}
Potrzebujemy wykryć i opisać zbiór konfiguracji, do których sterownik nie powinien dopuszczać.
\end{block}
\vspace{0.2em}
\begin{center}
\begin{tikzpicture}[scale=0.82]
 \draw[thick,fill=blue!5] (-1.1,0.3) rectangle (1.1,0.75);
 \foreach \x/\dx in {-0.8/-0.85,0.8/0.85} {
   \draw[ultra thick] (\x,0.3)--(\x+\dx,-0.75);
   \draw[ultra thick] (\x+\dx,-0.75)--(\x+1.25*\dx,-1.35);
   \fill[warn] (\x+1.25*\dx,-1.40) circle (0.08);
 }
 \draw[black!40,thick] (-2.3,-1.5)--(2.3,-1.5);
 \node[font=\scriptsize,align=center] at (0,-1.85) {duży rozstaw + tarcie = trudny powrót};
\end{tikzpicture}
\end{center}
\end{columns}
\note{Slajd ma przełożyć intuicję użytkownika na problem sterowania: niepożądana konfiguracja może wynikać z połączenia geometrii, tarcia, momentu i braku planowania odzysku.}
\end{frame}
